\documentclass[10pt]{article}

% ---------- Document Identity (update these only) ----------
\newcommand{\DocStage}{}
\newcommand{\DocTitle}{Your Road to Tier One SOF}
\newcommand{\ProgramName}{182-Week Civilian-to-Selection Readiness Pipeline}
\newcommand{\ProgramSubtitle}{}

% ---------- Page ----------
\usepackage[legalpaper,left=0.5in,right=0.5in,top=0.6in,bottom=0.45in]{geometry}

% ---------- Typography ----------
\usepackage[T1]{fontenc}
\usepackage[utf8]{inputenc}
\usepackage{libertinus}
\usepackage{microtype}
\usepackage{parskip}
\usepackage{ragged2e}
\usepackage{textcomp}
\AtBeginDocument{\color{black}}
\AtBeginDocument{\pagecolor{white}}
\AtBeginDocument{\justifying}

% ---------- Landscape pages ----------
\usepackage{pdflscape}

% ---------- Color ----------
\usepackage[table]{xcolor}
\definecolor{StageBlue}{HTML}{1F4E79}
\definecolor{StageBlueDark}{HTML}{163A5A}
\definecolor{LightGray}{HTML}{F2F4F7}
\definecolor{MidGray}{HTML}{D9DEE5}
\definecolor{RuleGray}{HTML}{C7CED8}
\definecolor{HeaderInk}{HTML}{000000}
\definecolor{Ink}{HTML}{000000}

% ---------- Utility ----------
\usepackage{enumitem}
\sloppy
\setlength{\emergencystretch}{2em}

% ---------- Boxes / UI ----------
\usepackage[most]{tcolorbox}

\tcbset{
  charterTitle/.style={
    enhanced,
    colback=StageBlue!4,
    colframe=StageBlue,
    boxrule=1.25pt,
    arc=3pt,
    left=16pt,right=16pt,top=14pt,bottom=14pt,
    borderline north={3.2pt}{0pt}{StageBlue},
    borderline west={4pt}{0pt}{StageBlue},
    borderline south={1.25pt}{0pt}{StageBlue},
    drop shadow={black!25!white},
  },
  charterRule/.style={
    enhanced,
    colback=LightGray,
    colframe=RuleGray,
    boxrule=0.85pt,
    arc=3pt,
    left=12pt,right=12pt,top=10pt,bottom=10pt,
    borderline west={3pt}{0pt}{StageBlue},
  },
  charterBadge/.style={
    enhanced,
    colback=StageBlue,
    coltext=white,
    colframe=StageBlueDark,
    boxrule=0.6pt,
    arc=2pt,
    left=6pt,right=6pt,top=3pt,bottom=3pt,
  }
}
\newtcbox{\Badge}{nobeforeafter,charterBadge}

% ---------- Lists ----------
\setlist[itemize]{leftmargin=1.5em, itemsep=0.25em, topsep=0.25em}
\setlist[enumerate]{leftmargin=1.8em, itemsep=0.25em, topsep=0.25em}

% ---------- TOC ----------
\usepackage{tocloft}
\setcounter{tocdepth}{3}
\setlength{\cftbeforesecskip}{3pt}
\setlength{\cftbeforesubsecskip}{2pt}
\setlength{\cftbeforesubsubsecskip}{0pt}
\renewcommand{\cftsecfont}{\bfseries}
\renewcommand{\cftsecpagefont}{\bfseries}
\renewcommand{\cftsubsecfont}{\normalfont}
\renewcommand{\cftsubsecpagefont}{\normalfont}
\renewcommand{\cftsubsubsecfont}{\normalfont}
\renewcommand{\cftsubsubsecpagefont}{\normalfont}
\setlength{\cftsecindent}{0pt}
\setlength{\cftsubsecindent}{1.25em}
\setlength{\cftsubsubsecindent}{2.8em}
\setlength{\cftsecnumwidth}{2.1em}
\setlength{\cftsubsecnumwidth}{2.8em}
\setlength{\cftsubsubsecnumwidth}{3.6em}

% Remove the built-in TOC heading so only the custom heading is shown
\makeatletter
\renewcommand{\tableofcontents}{\@starttoc{toc}}
\makeatother

% ---------- Header/Footer: page number only ----------
\usepackage{fancyhdr}
\pagestyle{fancy}
\fancyhf{}
\renewcommand{\headrulewidth}{0pt}
\renewcommand{\footrulewidth}{0pt}
\fancyfoot[C]{\thepage}

% ---------- Section formatting ----------
\usepackage{titlesec}

\titleformat{\section}
  {\Large\bfseries\color{Ink}}
  {\thesection}{0.6em}{}
  [\vspace{0.25em}\textcolor{StageBlue}{\rule{\linewidth}{0.8pt}}]

\titleformat{\subsection}
  {\large\bfseries\color{Ink}}
  {\thesubsection}{0.6em}{}

\titleformat{\subsubsection}
  {\normalsize\bfseries\color{Ink}}
  {\thesubsubsection}{0.6em}{}

\titlespacing*{\section}{0pt}{0.95em}{0.35em}
\titlespacing*{\subsection}{0pt}{0.65em}{0.25em}
\titlespacing*{\subsubsection}{0pt}{0.55em}{0.2em}

% ---------- Table engine ----------
\usepackage{tabularray}
\UseTblrLibrary{booktabs}

% ---------- tabularray: GLOBAL table treatment ----------
\SetTblrInner{
  cells = {fg=black, bg=white, valign=t},
}

% ---------- Header helpers ----------
\newcommand{\Hdr}[1]{\shortstack[c]{\strut #1\strut}}

% ---------- Lines ----------
\newcommand{\TitleLine}{\textcolor{StageBlue}{\rule{0.98\linewidth}{0.95pt}}}
\newcommand{\SubTitleLine}{\textcolor{RuleGray}{\rule{0.98\linewidth}{0.65pt}}}

% ---------- Continued on next page ----------
\newcommand{\ContinuedOnNextPageInline}{%
  \par
  \vfill
  \noindent\textcolor{RuleGray}{\rule{\linewidth}{1pt}}\vspace{-2pt}
  \noindent\hfill\textit{Continued on next page}\par\vspace{-10pt}
  \noindent\textcolor{RuleGray}{\rule{\linewidth}{1pt}}
  \clearpage
}

% ---------- PDF hyperlinks ----------
\usepackage[hidelinks]{hyperref}
\hypersetup{
  colorlinks=false,
  pdfborder={0 0 0},
  linktoc=all,
  pdftitle={\DocTitle},
  pdfauthor={\ProgramName}
}

% =========================================================
\begin{document}

\section*{Stage 6.D — Minimal vs Ideal Kit Pathways}
\phantomsection
\addcontentsline{toc}{section}{Stage 6.D — Minimal vs Ideal Kit Pathways}

\subsection*{6.D.1 Narrative}
\phantomsection
\addcontentsline{toc}{subsection}{6.D.1 Narrative}

\subsubsection*{Minimal Path definition}
\phantomsection

Minimal Path is the minimum deployable capability pathway that keeps every phase executable under start-from-zero constraints while preserving safety containment, validity prerequisites for admissible evidence, and continuity under real-world friction (weather, access, fatigue drift). Minimal Path aligns to the baseline Tier timeline in Stage 6.B. For any given phase, Minimal Path means: every category that is Tier 1 or higher in Stage 6.B for that phase is present, stable, and usable by the time gate-grade evidence is required.

\subsubsection*{Ideal Path definition}
\phantomsection

Ideal Path is a governed pull-forward overlay that increases reliability and reduces variance without changing what is “required.” Ideal Path selectively pulls forward Tier 2–Tier 3 capabilities only when they materially reduce (a) missed sessions due to access or weather friction, (b) invalid or non-comparable test events due to tool drift, and (c) predictable injury risk due to foot, load, or recovery-interface failures. Ideal Path is not “max spending.” It is risk-managed investment that targets the failure modes most likely to break gate windows.

\subsubsection*{Non-negotiable governing rule for phase execution}
\phantomsection

No phase may be executed without its Minimal Path Tier 1 prerequisites across applicable categories. If Minimal Path readiness is not met, the correct posture is delay escalation, reslot gate attempts inside the next protected window, and default \textbf{HOLD} for phase advancement. Do not improvise unsafe substitutions to “force” a timeline.

\subsubsection*{Validity substitution constraint tied to Stage 2 admissibility}
\phantomsection

Substitutions are permitted only when they preserve (1) safety boundaries, (2) Stage 2 comparability requirements for the specific measurement construct, and (3) Stage 1.F admissibility posture (\textbf{VALID} evidence only for gates). If a substitution changes the measurement construct (different tool class, different environment unit class, different route/pool verification posture) and comparability cannot be preserved, the resulting evidence is not gate-admissible and must be reslotted rather than “accepted anyway.”

\subsubsection*{Aquatics governance note}
\phantomsection

Aquatics capability (7.17) supports surface swim feasibility and safety logistics only. Underwater/breath-hold remains governance-gated and supervised-only; it must never be forced by equipment ownership. If governance, access, or supervision cannot be guaranteed, underwater remains “not scheduled; inadmissible without governance” regardless of Ideal Path status.

\subsubsection*{Freeze-window discipline for validity protection}
\phantomsection

Treat shoes, ruck interfaces, timing devices, and supplements as system variables. Do not introduce first-time or changed interfaces inside Deload+Test windows or during gate attempts. Stabilize any change in earlier work weeks so protected windows remain validity-clean and low variance.

\clearpage
\begin{landscape}

\subsection*{6.D.2 Summary Table}
\phantomsection
\addcontentsline{toc}{subsection}{6.D.2 Summary Table}

\begingroup
\scriptsize
\setlength{\tabcolsep}{2pt}
\renewcommand{\arraystretch}{1.12}

\begin{longtblr}{
  width=\linewidth,
  rowhead=1,
  colspec = {
    Q[l,wd=0.42in]
    X[l,1.85]
    X[l,3.20]
    X[l,3.20]
    X[l,3.40]
  },
  hlines,
  vlines,
  cells = {hsep=2pt, vsep=6pt, halign=l, valign=t},
  row{odd} = {bg=white},
  row{even} = {bg=LightGray},
  row{1} = {bg=StageBlue!15, fg=black, font=\bfseries, halign=l, valign=m},
}
\Hdr{Phase \#} &
\Hdr{Phase Name} &
\Hdr{Minimal Path Key Tier Levels by Category Group} &
\Hdr{Ideal Path Key Tier Levels by Category Group} &
\Hdr{Comments} \\

00 &
Phase 00 — Bodily Reactivation and Baseline Creation &
\textbf{HW}: \textbf{STR}/\textbf{COND}/\textbf{RUCK} = T0/T0/T0\newline
\textbf{FR}: \textbf{FEET}/\textbf{REC}/\textbf{PROTECT} = T1/T1/T0\newline
\textbf{ML}: \textbf{BIO}/\textbf{LOG}/\textbf{TEST} = T0/T1/T1\newline
\textbf{NH}: \textbf{NUT}/\textbf{HYG} = T1/T1\newline
\textbf{AP}: \textbf{ENV}/\textbf{GEN} = T1/T1\newline
\textbf{SM}: T1\newline
\textbf{AQ}: T0 &
\textbf{HW}: \textbf{STR}/\textbf{COND}/\textbf{RUCK} = T1/T1/T0\newline
\textbf{FR}: \textbf{FEET}/\textbf{REC}/\textbf{PROTECT} = T1/T1/T1\newline
\textbf{ML}: \textbf{BIO}/\textbf{LOG}/\textbf{TEST} = T0/T1/T2\newline
\textbf{NH}: \textbf{NUT}/\textbf{HYG} = T1/T1\newline
\textbf{AP}: \textbf{ENV}/\textbf{GEN} = T1/T1\newline
\textbf{SM}: T2\newline
\textbf{AQ}: T0 &
Minimal Path is strict “start-from-zero essentials” for admissibility and safe continuity. Ideal pull-forward highest ROI: indoor \textbf{COND} continuity (weather-proof), early \textbf{TEST} tool redundancy (avoid invalid logs), basic \textbf{PROTECT} to reduce predictable foot/skin failures. Freeze-window discipline is decisive even here: avoid changing shoes, timer, and log system near Weeks 4, 8, 12, and 16 protected windows. \\

01 &
Phase 01 — Base Conditioning and Hypertrophy &
\textbf{HW}: \textbf{STR}/\textbf{COND}/\textbf{RUCK} = T1/T1/T0\newline
\textbf{FR}: \textbf{FEET}/\textbf{REC}/\textbf{PROTECT} = T1/T1/T0\newline
\textbf{ML}: \textbf{BIO}/\textbf{LOG}/\textbf{TEST} = T0/T1/T1\newline
\textbf{NH}: \textbf{NUT}/\textbf{HYG} = T1/T1\newline
\textbf{AP}: \textbf{ENV}/\textbf{GEN} = T1/T1\newline
\textbf{SM}: T1\newline
\textbf{AQ}: T0 &
\textbf{HW}: \textbf{STR}/\textbf{COND}/\textbf{RUCK} = T1/T2/T0\newline
\textbf{FR}: \textbf{FEET}/\textbf{REC}/\textbf{PROTECT} = T1/T2/T1\newline
\textbf{ML}: \textbf{BIO}/\textbf{LOG}/\textbf{TEST} = T1/T2/T2\newline
\textbf{NH}: \textbf{NUT}/\textbf{HYG} = T1/T1\newline
\textbf{AP}: \textbf{ENV}/\textbf{GEN} = T1/T1\newline
\textbf{SM}: T2\newline
\textbf{AQ}: T0 &
Minimal Path introduces minimum deployable \textbf{STR} and \textbf{COND} capability for consistent execution. Ideal reduces missed sessions and invalid tests: pull \textbf{TEST} tool robustness earlier and add \textbf{REC}/\textbf{MOB} consistency to reduce overuse drift as volume rises. \\

02 &
Phase 02 — Maximal Strength Development &
\textbf{HW}: \textbf{STR}/\textbf{COND}/\textbf{RUCK} = T1/T1/T0\newline
\textbf{FR}: \textbf{FEET}/\textbf{REC}/\textbf{PROTECT} = T1/T1/T0\newline
\textbf{ML}: \textbf{BIO}/\textbf{LOG}/\textbf{TEST} = T0/T1/T1\newline
\textbf{NH}: \textbf{NUT}/\textbf{HYG} = T1/T1\newline
\textbf{AP}: \textbf{ENV}/\textbf{GEN} = T1/T1\newline
\textbf{SM}: T1\newline
\textbf{AQ}: T0 &
\textbf{HW}: \textbf{STR}/\textbf{COND}/\textbf{RUCK} = T2/T2/T0\newline
\textbf{FR}: \textbf{FEET}/\textbf{REC}/\textbf{PROTECT} = T2/T2/T1\newline
\textbf{ML}: \textbf{BIO}/\textbf{LOG}/\textbf{TEST} = T1/T2/T2\newline
\textbf{NH}: \textbf{NUT}/\textbf{HYG} = T1/T1\newline
\textbf{AP}: \textbf{ENV}/\textbf{GEN} = T1/T1\newline
\textbf{SM}: T2\newline
\textbf{AQ}: T0 &
Ideal pull-forward targets reliability of \textbf{STR} exposure and measurement comparability (\textbf{STR} implement stability, \textbf{TEST} tool stability). Freeze-window discipline: do not change \textbf{STR} implement class or logging method inside protected windows if tests depend on comparability. \\

03 &
Phase 03 — Converting Strength to Power and Endurance for SOF Selection &
\textbf{HW}: \textbf{STR}/\textbf{COND}/\textbf{RUCK} = T1/T1/T1\newline
\textbf{FR}: \textbf{FEET}/\textbf{REC}/\textbf{PROTECT} = T1/T1/T1\newline
\textbf{ML}: \textbf{BIO}/\textbf{LOG}/\textbf{TEST} = T0/T1/T2\newline
\textbf{NH}: \textbf{NUT}/\textbf{HYG} = T1/T1\newline
\textbf{AP}: \textbf{ENV}/\textbf{GEN} = T1/T1\newline
\textbf{SM}: T1\newline
\textbf{AQ}: T0 &
\textbf{HW}: \textbf{STR}/\textbf{COND}/\textbf{RUCK} = T2/T2/T2\newline
\textbf{FR}: \textbf{FEET}/\textbf{REC}/\textbf{PROTECT} = T2/T2/T2\newline
\textbf{ML}: \textbf{BIO}/\textbf{LOG}/\textbf{TEST} = T1/T2/T2\newline
\textbf{NH}: \textbf{NUT}/\textbf{HYG} = T1/T1\newline
\textbf{AP}: \textbf{ENV}/\textbf{GEN} = T1/T1\newline
\textbf{SM}: T2\newline
\textbf{AQ}: T0 &
Minimal Path introduces first minimum deployable \textbf{RUCK} capability by Phase 03 end gate. Ideal pull-forward reduces ruck-interface novelty risk: stabilize fit/config well before protected windows and improve \textbf{PROTECT} + \textbf{REC} to prevent hotspots and gait drift. \\

04 &
Phase 04 — Strength-Endurance and Work Capacity Development &
\textbf{HW}: \textbf{STR}/\textbf{COND}/\textbf{RUCK} = T1/T1/T1\newline
\textbf{FR}: \textbf{FEET}/\textbf{REC}/\textbf{PROTECT} = T1/T1/T1\newline
\textbf{ML}: \textbf{BIO}/\textbf{LOG}/\textbf{TEST} = T0/T1/T2\newline
\textbf{NH}: \textbf{NUT}/\textbf{HYG} = T1/T1\newline
\textbf{AP}: \textbf{ENV}/\textbf{GEN} = T1/T1\newline
\textbf{SM}: T1\newline
\textbf{AQ}: T1 &
\textbf{HW}: \textbf{STR}/\textbf{COND}/\textbf{RUCK} = T2/T2/T2\newline
\textbf{FR}: \textbf{FEET}/\textbf{REC}/\textbf{PROTECT} = T2/T2/T2\newline
\textbf{ML}: \textbf{BIO}/\textbf{LOG}/\textbf{TEST} = T1/T2/T2\newline
\textbf{NH}: \textbf{NUT}/\textbf{HYG} = T1/T1\newline
\textbf{AP}: \textbf{ENV}/\textbf{GEN} = T1/T1\newline
\textbf{SM}: T2\newline
\textbf{AQ}: T1 &
Minimal Path introduces first minimum deployable \textbf{AQUA} capability (surface-only feasibility and safety logistics). Ideal focuses on reducing invalid swim evidence (pool verification/logging fidelity) while preserving governance: underwater remains conditional and never forced. \\

05 &
Phase 05 — Aerobic Base and Extended Stamina Development &
\textbf{HW}: \textbf{STR}/\textbf{COND}/\textbf{RUCK} = T1/T1/T1\newline
\textbf{FR}: \textbf{FEET}/\textbf{REC}/\textbf{PROTECT} = T1/T1/T1\newline
\textbf{ML}: \textbf{BIO}/\textbf{LOG}/\textbf{TEST} = T0/T1/T2\newline
\textbf{NH}: \textbf{NUT}/\textbf{HYG} = T1/T1\newline
\textbf{AP}: \textbf{ENV}/\textbf{GEN} = T1/T1\newline
\textbf{SM}: T1\newline
\textbf{AQ}: T1 &
\textbf{HW}: \textbf{STR}/\textbf{COND}/\textbf{RUCK} = T2/T2/T2\newline
\textbf{FR}: \textbf{FEET}/\textbf{REC}/\textbf{PROTECT} = T2/T2/T2\newline
\textbf{ML}: \textbf{BIO}/\textbf{LOG}/\textbf{TEST} = T1/T2/T2\newline
\textbf{NH}: \textbf{NUT}/\textbf{HYG} = T2/T1\newline
\textbf{AP}: \textbf{ENV}/\textbf{GEN} = T2/T1\newline
\textbf{SM}: T2\newline
\textbf{AQ}: T1 &
Ideal pull-forward highest ROI: \textbf{ENV} apparel expansion and nutrition/hydration robustness to prevent environment-driven reslots and endurance-window invalidity. Freeze-window discipline becomes more consequential: route/tool/footwear changes near long events increase variance and threaten admissibility. \\

06 &
Phase 06 — Lactate Threshold and Power-Endurance Development &
\textbf{HW}: \textbf{STR}/\textbf{COND}/\textbf{RUCK} = T2/T2/T2\newline
\textbf{FR}: \textbf{FEET}/\textbf{REC}/\textbf{PROTECT} = T2/T2/T2\newline
\textbf{ML}: \textbf{BIO}/\textbf{LOG}/\textbf{TEST} = T1/T2/T2\newline
\textbf{NH}: \textbf{NUT}/\textbf{HYG} = T2/T1\newline
\textbf{AP}: \textbf{ENV}/\textbf{GEN} = T2/T1\newline
\textbf{SM}: T2\newline
\textbf{AQ}: T2 &
\textbf{HW}: \textbf{STR}/\textbf{COND}/\textbf{RUCK} = T2/T2/T2\newline
\textbf{FR}: \textbf{FEET}/\textbf{REC}/\textbf{PROTECT} = T2/T2/T2\newline
\textbf{ML}: \textbf{BIO}/\textbf{LOG}/\textbf{TEST} = T1/T2/T3\newline
\textbf{NH}: \textbf{NUT}/\textbf{HYG} = T2/T1\newline
\textbf{AP}: \textbf{ENV}/\textbf{GEN} = T2/T1\newline
\textbf{SM}: T2\newline
\textbf{AQ}: T2 &
Minimal Path reaches Tier 2 across most systems for continuity and validity. Ideal pull-forward is narrowly targeted: test-tool redundancy (avoid invalid gate weeks) and proactive \textbf{FEET}/\textbf{PROTECT} stability to prevent ruck/run interference and last-minute reslots. \\

07 &
Phase 07 — Advanced Hybrid Overload &
\textbf{HW}: \textbf{STR}/\textbf{COND}/\textbf{RUCK} = T2/T2/T2\newline
\textbf{FR}: \textbf{FEET}/\textbf{REC}/\textbf{PROTECT} = T2/T2/T2\newline
\textbf{ML}: \textbf{BIO}/\textbf{LOG}/\textbf{TEST} = T1/T2/T2\newline
\textbf{NH}: \textbf{NUT}/\textbf{HYG} = T2/T1\newline
\textbf{AP}: \textbf{ENV}/\textbf{GEN} = T2/T1\newline
\textbf{SM}: T2\newline
\textbf{AQ}: T2 &
\textbf{HW}: \textbf{STR}/\textbf{COND}/\textbf{RUCK} = T2/T2/T2\newline
\textbf{FR}: \textbf{FEET}/\textbf{REC}/\textbf{PROTECT} = T2/T2/T2\newline
\textbf{ML}: \textbf{BIO}/\textbf{LOG}/\textbf{TEST} = T1/T2/T3\newline
\textbf{NH}: \textbf{NUT}/\textbf{HYG} = T2/T1\newline
\textbf{AP}: \textbf{ENV}/\textbf{GEN} = T2/T1\newline
\textbf{SM}: T3\newline
\textbf{AQ}: T2 &
Hybrid overload increases drift risk: Ideal focuses on \textbf{SELF-MGMT} redundancy (planning discipline, adherence control) and \textbf{TEST} tool redundancy to reduce invalidity and stacking temptations. \\

08 &
Phase 08 — Structural Durability and Load Tolerance &
\textbf{HW}: \textbf{STR}/\textbf{COND}/\textbf{RUCK} = T2/T2/T2\newline
\textbf{FR}: \textbf{FEET}/\textbf{REC}/\textbf{PROTECT} = T2/T2/T2\newline
\textbf{ML}: \textbf{BIO}/\textbf{LOG}/\textbf{TEST} = T1/T2/T2\newline
\textbf{NH}: \textbf{NUT}/\textbf{HYG} = T2/T1\newline
\textbf{AP}: \textbf{ENV}/\textbf{GEN} = T2/T1\newline
\textbf{SM}: T2\newline
\textbf{AQ}: T2 &
\textbf{HW}: \textbf{STR}/\textbf{COND}/\textbf{RUCK} = T2/T2/T3\newline
\textbf{FR}: \textbf{FEET}/\textbf{REC}/\textbf{PROTECT} = T3/T2/T3\newline
\textbf{ML}: \textbf{BIO}/\textbf{LOG}/\textbf{TEST} = T1/T2/T3\newline
\textbf{NH}: \textbf{NUT}/\textbf{HYG} = T2/T1\newline
\textbf{AP}: \textbf{ENV}/\textbf{GEN} = T2/T1\newline
\textbf{SM}: T2\newline
\textbf{AQ}: T2 &
Ideal pull-forward emphasizes ruck and foot-interface redundancy (reduce blister/gait drift probability) and \textbf{PROTECT} posture to prevent predictable breakdown during long-load milestones. Freeze-window discipline is strict: no first-time ruck or boot interface changes inside protected windows. \\

09 &
Phase 09 — High-Volume Calisthenics and Stamina &
\textbf{HW}: \textbf{STR}/\textbf{COND}/\textbf{RUCK} = T2/T2/T2\newline
\textbf{FR}: \textbf{FEET}/\textbf{REC}/\textbf{PROTECT} = T2/T2/T2\newline
\textbf{ML}: \textbf{BIO}/\textbf{LOG}/\textbf{TEST} = T1/T2/T2\newline
\textbf{NH}: \textbf{NUT}/\textbf{HYG} = T2/T1\newline
\textbf{AP}: \textbf{ENV}/\textbf{GEN} = T2/T1\newline
\textbf{SM}: T2\newline
\textbf{AQ}: T2 &
\textbf{HW}: \textbf{STR}/\textbf{COND}/\textbf{RUCK} = T2/T2/T3\newline
\textbf{FR}: \textbf{FEET}/\textbf{REC}/\textbf{PROTECT} = T3/T3/T3\newline
\textbf{ML}: \textbf{BIO}/\textbf{LOG}/\textbf{TEST} = T1/T2/T3\newline
\textbf{NH}: \textbf{NUT}/\textbf{HYG} = T2/T1\newline
\textbf{AP}: \textbf{ENV}/\textbf{GEN} = T2/T1\newline
\textbf{SM}: T2\newline
\textbf{AQ}: T2 &
Ideal pull-forward concentrates on \textbf{REC}/\textbf{MOB} and \textbf{PROTECT} to reduce high-volume \textbf{CAL} overuse drift and preserve repeatable test windows. This is reliability-focused, not performance-driven. \\

10 &
Phase 10 — Advanced Tactical Readiness &
\textbf{HW}: \textbf{STR}/\textbf{COND}/\textbf{RUCK} = T2/T2/T2\newline
\textbf{FR}: \textbf{FEET}/\textbf{REC}/\textbf{PROTECT} = T2/T2/T2\newline
\textbf{ML}: \textbf{BIO}/\textbf{LOG}/\textbf{TEST} = T1/T2/T2\newline
\textbf{NH}: \textbf{NUT}/\textbf{HYG} = T2/T1\newline
\textbf{AP}: \textbf{ENV}/\textbf{GEN} = T2/T1\newline
\textbf{SM}: T2\newline
\textbf{AQ}: T2 &
\textbf{HW}: \textbf{STR}/\textbf{COND}/\textbf{RUCK} = T2/T2/T3\newline
\textbf{FR}: \textbf{FEET}/\textbf{REC}/\textbf{PROTECT} = T3/T3/T3\newline
\textbf{ML}: \textbf{BIO}/\textbf{LOG}/\textbf{TEST} = T1/T2/T3\newline
\textbf{NH}: \textbf{NUT}/\textbf{HYG} = T2/T1\newline
\textbf{AP}: \textbf{ENV}/\textbf{GEN} = T2/T1\newline
\textbf{SM}: T3\newline
\textbf{AQ}: T2 &
Ideal pull-forward targets the categories most likely to invalidate multi-domain gate days: \textbf{TEST} tools, \textbf{FEET}/\textbf{REC}/\textbf{PROTECT} stability, and \textbf{SELF-MGMT} to enforce reslot discipline instead of stacking. Aquatics remains surface-focused; underwater remains governance-gated and may remain not scheduled. \\

11 &
Phase 11 — Pre-Selection Conditioning and Recovery &
\textbf{HW}: \textbf{STR}/\textbf{COND}/\textbf{RUCK} = T2/T2/T2\newline
\textbf{FR}: \textbf{FEET}/\textbf{REC}/\textbf{PROTECT} = T2/T2/T2\newline
\textbf{ML}: \textbf{BIO}/\textbf{LOG}/\textbf{TEST} = T1/T2/T2\newline
\textbf{NH}: \textbf{NUT}/\textbf{HYG} = T2/T1\newline
\textbf{AP}: \textbf{ENV}/\textbf{GEN} = T2/T1\newline
\textbf{SM}: T2\newline
\textbf{AQ}: T2 &
\textbf{HW}: \textbf{STR}/\textbf{COND}/\textbf{RUCK} = T2/T2/T3\newline
\textbf{FR}: \textbf{FEET}/\textbf{REC}/\textbf{PROTECT} = T3/T3/T3\newline
\textbf{ML}: \textbf{BIO}/\textbf{LOG}/\textbf{TEST} = T1/T2/T3\newline
\textbf{NH}: \textbf{NUT}/\textbf{HYG} = T2/T1\newline
\textbf{AP}: \textbf{ENV}/\textbf{GEN} = T2/T1\newline
\textbf{SM}: T3\newline
\textbf{AQ}: T2 &
Ideal pull-forward continues the “variance minimization” posture: reduce late-phase failures driven by foot breakdown, recovery collapse, and tool drift. Maintain freeze-window discipline rigorously. \\

12 &
Phase 12 — Final Pre-Selection Conditioning &
\textbf{HW}: \textbf{STR}/\textbf{COND}/\textbf{RUCK} = T2/T2/T3\newline
\textbf{FR}: \textbf{FEET}/\textbf{REC}/\textbf{PROTECT} = T3/T3/T3\newline
\textbf{ML}: \textbf{BIO}/\textbf{LOG}/\textbf{TEST} = T1/T3/T3\newline
\textbf{NH}: \textbf{NUT}/\textbf{HYG} = T2/T1\newline
\textbf{AP}: \textbf{ENV}/\textbf{GEN} = T2/T1\newline
\textbf{SM}: T3\newline
\textbf{AQ}: T3 &
\textbf{HW}: \textbf{STR}/\textbf{COND}/\textbf{RUCK} = T2/T2/T3\newline
\textbf{FR}: \textbf{FEET}/\textbf{REC}/\textbf{PROTECT} = T3/T3/T3\newline
\textbf{ML}: \textbf{BIO}/\textbf{LOG}/\textbf{TEST} = T2/T3/T3\newline
\textbf{NH}: \textbf{NUT}/\textbf{HYG} = T2/T1\newline
\textbf{AP}: \textbf{ENV}/\textbf{GEN} = T2/T1\newline
\textbf{SM}: T3\newline
\textbf{AQ}: T3 &
Minimal Path introduces selective Tier 3 “verification reliability” posture (redundancy where failure would invalidate decisive windows). Ideal adds limited monitoring robustness (\textbf{BIO} to T2) without changing what is required for deployability. \\

13 &
Phase 13 — Full-Spectrum Readiness Consolidation &
\textbf{HW}: \textbf{STR}/\textbf{COND}/\textbf{RUCK} = T2/T2/T3\newline
\textbf{FR}: \textbf{FEET}/\textbf{REC}/\textbf{PROTECT} = T3/T3/T3\newline
\textbf{ML}: \textbf{BIO}/\textbf{LOG}/\textbf{TEST} = T1/T3/T3\newline
\textbf{NH}: \textbf{NUT}/\textbf{HYG} = T2/T1\newline
\textbf{AP}: \textbf{ENV}/\textbf{GEN} = T2/T1\newline
\textbf{SM}: T3\newline
\textbf{AQ}: T3 &
\textbf{HW}: \textbf{STR}/\textbf{COND}/\textbf{RUCK} = T2/T2/T3\newline
\textbf{FR}: \textbf{FEET}/\textbf{REC}/\textbf{PROTECT} = T3/T3/T3\newline
\textbf{ML}: \textbf{BIO}/\textbf{LOG}/\textbf{TEST} = T2/T3/T3\newline
\textbf{NH}: \textbf{NUT}/\textbf{HYG} = T2/T1\newline
\textbf{AP}: \textbf{ENV}/\textbf{GEN} = T2/T1\newline
\textbf{SM}: T3\newline
\textbf{AQ}: T3 &
Consolidation maintains Tier 3 stability posture and enforces “no novelty near protected windows” as a hard validity control. Ideal differs mainly by earlier monitoring robustness; neither changes what is required for deployability. \\

\end{longtblr}

\endgroup

\end{landscape}

\end{document}